% The mixture factor graph of Section~\ref{sec:mixture}, drawn by hand. Input
% by textbook.tex; a sketch and not a rendered figure, so it is outside the
% figure cache and the render cap, which govern generated output only.
\begin{figure}[htbp]
\centering
\begin{tikzpicture}[
    scale=1.0
]
  \foreach \i in {1,...,5} {
    \node[fgvariable] (z\i) at (1.7*\i, 1.6) {$z_{\i}$};
    \node[fgdata] (e\i) at (1.7*\i, 0.8) {};
    \node[fgobserved] (y\i) at (1.7*\i, 0.1) {$y_{\i}$};
    \draw[fglink] (z\i) -- (e\i) -- (y\i);
  }
  % Where the chain would be, had it not been removed.
  \foreach \i/\next in {1/2, 2/3, 3/4, 4/5} {
    \draw[fgabsent] (z\i) -- (z\next);
  }
  \node[fgnote] at (9.0, 1.6)
    {no pairwise factor};
  \node[fgnote] at (9.0, 0.8)
    {$\psi_i = w_{z_i}\,\mathcal{N}(y_i;\, \mu_{z_i}, \sigma_{z_i}^2)$};
\end{tikzpicture}
\caption{The Gaussian mixture as a factor graph: $N$ component labels $z_i$
over $K$ values, one shaded unary factor each folding in one observation
$y_i$, and no pairwise factor at all. The dashed lines mark where the
transition factors of Figure~\ref{fig:hmm:sketch} would sit; removing them is
exactly what turns the hidden Markov model into the mixture, and it turns
Equation~\ref{eq:sum-product} into the per-variable normalization that is
Equation~\ref{eq:responsibilities}. The same emission family serves both, and
receives the same responsibilities from either. A sketch of the structure and
not a rendering of any instance.}
\label{fig:mixture:sketch}
\end{figure}
